% Image credits for the photographs and AI illustrations of Book 1 (Primary and Middle School), Indonesian edition.
% Machine-readable license records live in images/book1/CREDITS.md; keep this
% file in sync with frontmatter/image-credits.tex (one \item per item there).
% Names, institutions, image identifiers and licences are kept exactly as in
% the English file; only the frame and the descriptions are translated.
% "Viñas" is written in UTF-8, as the English file now writes it too.
\chapter*{Kredit Gambar}

Gambar, peta, dan grafik dalam buku ini adalah karya asli; peta dan grafik
yang digambar dari data yang telah diterbitkan menyebutkan sumbernya di
keterangan gambar atau di daftar ini. Foto-foto dipakai di bawah lisensi
bebasnya masing-masing, dengan terima kasih kepada para pembuatnya; tautan
sumber lengkap dan alamat lisensi setiap foto tercantum dalam
\texttt{images/book1/CREDITS.md} di repositori buku ini.

\bigskip
Selain foto dan gambar asli, sebagian gambar adalah ilustrasi buatan AI,
yang dihasilkan dengan pembangkit gambar OpenAI dari perintah yang ditulis
oleh para penyunting buku ini, dan diperiksa ketepatan ilmiahnya sebelum
dimuat. Ilustrasi itu hanya menampilkan adegan umum, tidak pernah tempat
yang nyata, orang yang nyata, atau spesimen yang nyata. Perintah lengkap
untuk setiap gambar buatan AI dicatat dalam
\texttt{images/book1/ai/PROMPTS.md} di repositori.

\bigskip
\noindent Foto dan data, menurut bab:
\begin{itemize}\small
  \item Bab~\ref{ch:g2:water-on-earth} (air di Bumi): Bumi dilihat dari Apollo~17, 7~Desember 1972 (``The Blue Marble''); NASA, domain publik (melalui Wikimedia Commons). Ke-100 persegi dihitung dari relief NOAA ETOPO1; ke-100 ember dari tabel Water Science School milik U.S. Geological Survey (menurut Shiklomanov, 1993).
  \item Bab~\ref{ch:g3:stones-pebbles-sand} (batu, kerakal, dan pasir): sebuah dinding batu granit tanpa adukan; foto karya Tony Atkin, Geograph, CC BY-SA 2.0 (melalui Wikimedia Commons). Atap dari lempengan batu di Swiss; foto karya Michael C. Rygel, CC BY-SA 3.0 (melalui Wikimedia Commons).
  \item Bab~\ref{ch:g3:continents-and-oceans} (bola dunia): kedua bola dunia dan peta dunia digambar dari garis pantai Natural Earth (domain publik).
  \item Bab~\ref{ch:g4:clouds-rain-snow} (awan, hujan, dan salju): kumulus di atas Kansas (S.~Corfidi), petak-petak sirus (R.~F.~Kresge), selapis stratus di sekeliling puncak-puncak gunung, badai guntur di atas gurun Arizona (koleksi B.~Hooke), sebutir batu es yang sangat besar; semuanya NOAA Photo Library, domain publik (melalui Wikimedia Commons).
  \item Bab~\ref{ch:g4:volcanoes-earthquakes-first-look} (gunung api dan gempa bumi, pengenalan pertama): lava yang memasuki samudra di Kilauea, USGS Hawaiian Volcano Observatory; Gunung St.~Helens, 18~Mei 1980, foto karya Austin Post, USGS; sebuah rumah yang terlempar dari fondasinya oleh gempa bumi Loma Prieta tahun 1989, foto karya J.~K.~Nakata, USGS; semuanya domain publik (melalui Wikimedia Commons). Petanya memakai daftar gunung api Holosen dari Smithsonian Global Volcanism Program dan katalog gempa bumi U.S. Geological Survey.
  \item Bab~\ref{ch:g4:world-physical-map} (peta fisik): Himalaya dan Amazon yang dipotret dari Stasiun Antariksa Internasional, NASA, domain publik (melalui Wikimedia Commons). Peta fisiknya digambar dari relief NOAA ETOPO1, dengan sungai, danau, dan garis pantai dari Natural Earth (domain publik).
  \item Bab~\ref{ch:g5:water-shapes-the-land} (air membentuk daratan): Sungai Mamoré yang dipotret dari Stasiun Antariksa Internasional (ISS007-E-10797) dan delta Sungai Nil yang dilihat oleh instrumen MODIS pada satelit Terra, NASA, domain publik (melalui Wikimedia Commons).
  \item Bab~\ref{ch:g5:weather-and-climate} (cuaca dan iklim): Gurun Sonora, U.S. National Park Service, Saguaro National Park, domain publik (melalui Wikimedia Commons). Diagram iklimnya digambar dari NOAA NCEI U.S. Climate Normals 1991--2020 (domain publik).
  \item Bab~\ref{ch:g5:natural-hazards} (bahaya alam): mata Badai Isabel yang dipotret dari Stasiun Antariksa Internasional oleh Ed Lu, NASA, domain publik; tanah longsor tahun 2014 di negara bagian Washington, Mark Reid, U.S. Geological Survey, domain publik (keduanya melalui Wikimedia Commons).
  \item Bab~\ref{ch:g6:earth-a-planet} (Bumi, sebuah planet): Venus dari Mariner~10 (NASA/JPL-Caltech); Bumi dari Apollo~17 (NASA); \emph{Earthrise} dari Apollo~8 (NASA, Bill Anders); Mars dan Jupiter dari Teleskop Antariksa Hubble (NASA, ESA, Hubble Heritage Team, dan A.~Simon); semuanya domain publik (melalui Wikimedia Commons). \emph{Pale Blue Dot} yang dilihat Voyager~1 pada tahun 1990, versi 2020 (NASA/JPL-Caltech, PIA23645, domain publik). Data planet dari lembar fakta planet NASA.
  \item Bab~\ref{ch:g6:latitude-longitude} (lintang dan bujur): H4 karya Harrison di National Maritime Museum, Greenwich, foto karya Colonel Warden, CC BY-SA 3.0 (melalui Wikimedia Commons). Garis pantai dan kota dari Natural Earth (domain publik).
  \item Bab~\ref{ch:g6:rocks-and-minerals} (batuan dan mineral): granit, basal, batu gamping berfosil, batu pasir, marmer, serpih, sebuah kristal kuarsa, dan selembar mika muskovit, foto karya James St.~John, CC BY 2.0; sebuah atap batu sabak, foto karya Øyvind Holmstad, CC0 (semuanya melalui Wikimedia Commons).
  \item Bab~\ref{ch:g6:inside-the-earth} (bagian dalam Bumi): Inge Lehmann pada tahun 1932, foto karya Even Neuhaus, Royal Library of Denmark, domain publik (melalui Wikimedia Commons). Lapisan, massa jenis, dan kecepatan gelombangnya digambar dari Preliminary Reference Earth Model (A.~M. Dziewonski dan D.~L. Anderson, 1981).
  \item Bab~\ref{ch:g6:water-underground} (air di bawah tanah): Big Spring, Missouri, foto karya Johnston9494, CC BY-SA 4.0 (melalui Wikimedia Commons). Porositas menurut R.~C. Heath, \emph{Basic Ground-Water Hydrology}, U.S. Geological Survey (1983); bagian-bagian air tawar dari USGS Water Science School.
  \item Bab~\ref{ch:g7:plate-tectonics} (tektonika lempeng): benua-benua selatan yang disatukan kembali, dengan jalur-jalur fosil \emph{Cynognathus}, \emph{Lystrosaurus}, \emph{Mesosaurus}, dan \emph{Glossopteris}; U.S. Geological Survey, \emph{This Dynamic Earth}, domain publik.
  \item Bab~\ref{ch:g7:plate-tectonics}: sesar San Andreas yang melintasi Dataran Carrizo; foto karya Robert E. Wallace, U.S. Geological Survey, domain publik.
  \item Bab~\ref{ch:g7:plate-tectonics}: batas-batas lempeng menurut P.~Bird (2003), GeoJSON oleh H.~Ahlenius (ODC-By 1.0); garis pantai dari Natural Earth (domain publik); gempa bumi dari katalog U.S. Geological Survey; gunung api dari Smithsonian Global Volcanism Program.
  \item Bab~\ref{ch:g7:volcanoes} (gunung api): Gunung St.~Helens dari Johnston Ridge pada 17~Mei dan 10~September 1980, foto karya Harry Glicken, U.S. Geological Survey; sebuah air mancur lava di Kilauea, 1983, dan Mauna Loa dari udara, 1985, foto karya J.~D. Griggs, U.S. Geological Survey; tiang letusan Gunung Pinatubo pada 12~Juni 1991, foto karya Dave Harlow, U.S. Geological Survey; semuanya domain publik (melalui Wikimedia Commons). Peta dan hitungannya memakai daftar gunung api Holosen dari Smithsonian Global Volcanism Program; batas-batas lempeng menurut P.~Bird (2003).
  \item Bab~\ref{ch:g7:earthquakes} (gempa bumi): sebuah desa di dekat pesisir Sumatra setelah tsunami 26~Desember 2004, foto U.S. Navy karya Philip A. McDaniel; gawir-gawir sesar gempa bumi Borah Peak tahun 1983, foto karya R.~E. Wallace, U.S. Geological Survey; jalan yang amblas ke Danau Hebgen, 1959, foto karya I.~J. Witkind, U.S. Geological Survey; semuanya domain publik (melalui Wikimedia Commons). Hitungan dan kedalamannya memakai katalog gempa bumi U.S. Geological Survey; tabel intensitasnya diringkas dari skala Modified Mercalli milik U.S. Geological Survey.
  \item Bab~\ref{ch:g7:sediments-and-fossils} (sedimen dan fosil): Grand Canyon dari Hermits Rest, foto NPS karya Kristen M. Caldon, domain publik; sebuah amonit, sebuah trilobit, dan riak gelombang di Capitol Reef National Park, foto karya James St.~John, CC BY 2.0; jejak kaki dinosaurus di Dinosaur Ridge, foto karya Dexter Perkins, CC0 (semuanya melalui Wikimedia Commons).
  \item Bab~\ref{ch:g7:geological-time} (umur Bumi): James Hutton, potret karya Henry Raeburn, domain publik (melalui Wikimedia Commons). Umur dan warna skala waktunya adalah umur dan warna International Chronostratigraphic Chart (International Commission on Stratigraphy).
  \item Bab~\ref{ch:g8:the-atmosphere} (atmosfer): peluncuran balon cuaca di Topeka, Kansas, foto karya Stephen Corfidi, NOAA/National Weather Service; tepi Bumi saat matahari terbenam dari Stasiun Antariksa Internasional, NASA, ISS Expedition~23; keduanya domain publik (melalui Wikimedia Commons). Profil-profilnya dihitung dari \emph{U.S. Standard Atmosphere, 1976} (NOAA, NASA, dan U.S. Air Force).
  \item Bab~\ref{ch:g8:winds-global-circulation} (angin dan sirkulasi global): benua Amerika yang dilihat oleh satelit GOES-16 pada 26~Juli 2022, NOAA, domain publik (melalui Wikimedia Commons). Garis pantai pada peta angin pasat dari Natural Earth (domain publik).
  \item Bab~\ref{ch:g8:fronts-and-weather-maps} (front dan peta cuaca): sebuah depresi di lepas pantai Islandia yang dilihat oleh satelit Terra, NASA/GSFC, MODIS Rapid Response Team, Jacques Descloitres, domain publik (melalui Wikimedia Commons). Garis pantai pada peta massa udara dari Natural Earth (domain publik).
  \item Bab~\ref{ch:g8:the-oceans} (samudra): HMS \emph{Challenger} yang bersiap melakukan pengukuran kedalaman, 1872, gravur dari laporan-laporan ekspedisi Challenger yang direproduksi dalam W.~A. Herdman, \emph{Founders of Oceanography} (1923); Arus Teluk yang dilihat oleh instrumen MODIS pada satelit Aqua, NASA Earth Observatory; keduanya domain publik (melalui Wikimedia Commons). Profil dasar samudranya digambar dari relief NOAA ETOPO1; garis pantai pada peta pilin dan peta pelampung dari Natural Earth (domain publik).
  \item Bab~\ref{ch:g8:ice-on-earth} (es di Bumi): Gletser Pedersen pada tahun 1990 dan 2016, foto karya M.~Tetreau dan D.~Kurtz, U.S. National Park Service; tepi sebuah paparan es di lepas Semenanjung Antarktika, NASA, Operation IceBridge (M.-J. Viñas); semuanya domain publik (melalui Wikimedia Commons). Grafik es lautnya memakai NSIDC Sea Ice Index, versi~4 (F.~Fetterer dan kawan-kawan, 2025, doi:10.7265/a98x-0f50).
  \item Bab~\ref{ch:g9:climate-past-and-present} (iklim Bumi, dahulu dan sekarang): sepotong inti es GISP2 dari Greenland, NOAA Paleoclimatology, domain publik (melalui Wikimedia Commons). Rekaman-rekamannya digambar dari NASA GISTEMP versi~4, rekaman karbon dioksida Mauna Loa (NOAA Global Monitoring Laboratory dan Scripps Institution of Oceanography), suhu EPICA Dome~C (J.~Jouzel dan kawan-kawan, 2007), gabungan karbon dioksida inti es Antarktika (B.~Bereiter dan kawan-kawan, 2015), dan muka laut dari NOAA Laboratory for Satellite Altimetry; pernyataan-pernyataan penilaiannya adalah pernyataan IPCC Sixth Assessment Report, Working Group~I, Summary for Policymakers (2021).
  \item Bab~\ref{ch:g9:natural-risks} (hidup bersama risiko alam): peta bahaya Nevado del Ruiz, dari T.~L. Wright dan T.~C. Pierson, \emph{Living with Volcanoes}, U.S. Geological Survey Circular 1073 (1992); sebuah pelampung peringatan tsunami yang diperiksa oleh awak USCGC \emph{William Hart}, U.S. Coast Guard, 2021; keduanya domain publik (melalui Wikimedia Commons).
  \item Bab~\ref{ch:g9:earth-resources} (sumber daya dari Bumi): tambang Chuquicamata yang dilihat oleh ASTER, NASA/\allowbreak METI/\allowbreak AIST/\allowbreak Japan Space Systems dan U.S./\allowbreak Japan ASTER Science Team, 2007, domain publik; sebuah lapisan batu bara pada tebing potongan jalan di Belmont County, Ohio, foto karya James St.~John, CC BY 2.0 (keduanya melalui Wikimedia Commons). Faktor-faktor konsentrasinya dari D.~Perkins dan kawan-kawan, \emph{Mineralogy}, dan S.~Earle, \emph{Environmental Geology}.
  \item Bab~\ref{ch:g9:landforms} (bentuk lahan): Seven Sisters, foto karya Arild Vågen, 2017, CC BY-SA 4.0; gumuk-gumuk pasir berbentuk sabit di pesisir Brasil selatan, foto astronaut NASA dari Stasiun Antariksa Internasional, 2017, domain publik; stalaktit di Carlsbad Cavern, foto National Park Service karya Daniel Leifheit, 2018, domain publik (semuanya melalui Wikimedia Commons).
  \item Bab~\ref{ch:g9:maps-of-the-earth} (peta seluruh Bumi): delta Sungai Lena, citra Landsat tahun 2000, USGS EROS Data Center melalui NASA Visible Earth, domain publik; sebuah satelit GPS, gambaran seniman NASA, domain publik (keduanya melalui Wikimedia Commons). Garis pantai pada peta-peta dunia dari Natural Earth (domain publik).
\end{itemize}
